\documentclass[12pt,letterpaper, onecolumn]{exam}
\usepackage{amsmath}
\usepackage{amssymb}
\usepackage[lmargin=71pt, tmargin=1.2in]{geometry}
\usepackage{xcolor}
\usepackage[brazil]{babel}
\usepackage{natbib}
\usepackage{enumitem}
\usepackage{booktabs}
\usepackage[hidelinks]{hyperref}
\urlstyle{same}

\renewenvironment{solution}
{\par\noindent\textbf{R:}\enspace\color{blue}}
{\par\color{black}}

% \chead{\hline} % Un-comment to draw line below header
\thispagestyle{empty}

\begin{document}

\begingroup
    \centering
    \LARGE Lista 11 - Econometria I - 2026\\[0.5em]
    %\large \today\\[0.5em]
    \large Prof: André Luiz Pereira Mancha \par
    \large Monitor: Tuffy Licciardi Issa \par
    %\large Roll Number\par
    %\large Class/Batch/Section\par
\endgroup
\rule{\textwidth}{0.4pt}
\pointsdroppedatright
\printanswers
\renewcommand{\solutiontitle}{\noindent\textbf{Ans:}\enspace}
\newcommand{\qstar}{\hspace{-0.6em}\textsuperscript{*}\hspace{0.2em}}

\begin{questions}

\question (9.1, \citep{wooldridge2016}) No Exercício 4.11, o \(R\)-quadrado da estimativa do modelo
\[
\log(wage)=\beta_0+\beta_1\log(sales)+\beta_2\log(mktval)+\beta_3 profmarg+\beta_4 ceoten+\beta_5 comten+u,
\]
usando os dados contidos no arquivo \texttt{CEOSAL2}, era \(R^2=0{,}353\) \((n=177)\). Quando \(ceoten^2\) e \(comten^2\) são adicionados, \(R^2=0{,}375\). Existe evidência de má-especificação de forma funcional neste modelo?

\vspace{0.7cm}

\question\qstar (9.2, \citep{wooldridge2016}) Modifiquemos o Exercício em Computador C4, do Capítulo 8, utilizando os resultados das eleições de 1990 dos candidatos que foram eleitos em 1988. O Candidato A foi eleito em 1988 e buscava a reeleição em 1990; \textit{voteA90} é a porcentagem de votos do Candidato A na eleição de 1990. A porcentagem de votos do Candidato A na eleição de 1988 é usada como uma variável \textit{proxy} da qualidade do candidato. Todas as demais variáveis são das eleições de 1990. As seguintes equações foram estimadas, utilizando os dados contidos no arquivo \texttt{VOTE2}:

\[
\widehat{voteA90}=75{,}71+0{,}312\,prtystrA+4{,}93\,democA-0{,}929\,\log(expendA)-1{,}950\,\log(expendB)
\]
\[
\hspace{2.0cm}(9{,}25)\hspace{0.8cm}(0{,}046)\hspace{0.8cm}(1{,}01)\hspace{0.8cm}(0{,}684)\hspace{0.8cm}(0{,}281)
\]
\[
n=186,\; R^2=0{,}465
\]

e

\[
\begin{aligned}
\widehat{voteA90}=\;& 70{,}81+0{,}282\,prtystrA+4{,}52\,democA \\
& -0{,}839\,\log(expendA)-1{,}846\,\log(expendB)+0{,}067\,voteA88
\end{aligned}
\]
\[
\begin{aligned}
& (10{,}01)\hspace{0.8cm}(0{,}052)\hspace{0.8cm}(1{,}06) \\
& (0{,}687)\hspace{0.8cm}(0{,}292)\hspace{0.8cm}(0{,}053)
\end{aligned}
\]
\[
n=186,\; R^2=0{,}499,\; \bar{R}^2=0{,}485
\]

\begin{enumerate}[label=(\roman*)]
    \item Interprete o coeficiente de \textit{voteA88} e comente sobre sua significância estatística.
    \item A adição de \textit{voteA88} tem muito efeito sobre os outros coeficientes?
\end{enumerate}

\vspace{0.7cm}

\question\qstar (9.3, \citep{wooldridge2016}) Seja \textit{math10} a porcentagem de aprovação em um teste padrão de matemática de estudantes de uma escola secundária de Michigan (veja também o Exemplo 4.2). Estamos interessados em estimar o efeito do gasto por estudante no desempenho em matemática. Um modelo simples é
\[
math10=\beta_0+\beta_1\log(expend)+\beta_2\log(enroll)+\beta_3 poverty+u,
\]
em que \textit{poverty} é a porcentagem de estudantes vivendo em condições de pobreza.
\begin{enumerate}[label=(\roman*)]
    \item A variável \textit{lnchprg} é a porcentagem de estudantes qualificados para o programa de merenda escolar financiado pelo governo federal. Por que ela é uma variável \textit{proxy} razoável de \textit{poverty}?
    \item A tabela seguinte contém estimativas MQO, com e sem \textit{lnchprg} como uma variável explicativa:
\end{enumerate}

\begin{center}
\begin{tabular}{lcc}
\toprule
Variáveis independentes & (1) & (2) \\
\midrule
\(\log(expend)\) & 11,13 & 7,75 \\
& (3,30) & (3,04) \\
\(\log(enroll)\) & 0,022 & -1,26 \\
& (0,615) & (0,58) \\
\textit{lnchprg} & -- & -0,324 \\
& & (0,036) \\
intercepto & -69,24 & -23,14 \\
& (26,72) & (24,99) \\
Observações & 428 & 428 \\
\(R^2\) & 0,0297 & 0,1893 \\
\bottomrule
\end{tabular}
\end{center}

\begin{enumerate}[label=(\roman*),resume]
    \item Explique por que o efeito dos gastos (\textit{expend}) sobre \textit{math10} é menor na coluna (2) do que na coluna (1). O efeito na coluna (2) ainda é estatisticamente maior que zero?
    \item Parece que as taxas de aprovação são menores em escolas maiores, com os outros fatores sendo iguais? Explique.
    \item Interprete o coeficiente de \textit{lnchprg} na coluna (2).
    \item O que você deduz do substancial aumento de \(R^2\) da coluna (1) para a coluna (2)?
\end{enumerate}

\vspace{0.7cm}

\question (9.5, \citep{wooldridge2016}) No Exemplo 4.4, estimamos um modelo relacionando número de crimes no \textit{campus} às matrículas de estudantes em um grupo de faculdades. A amostra que usamos não era uma amostra aleatória de faculdades nos Estados Unidos, pois muitas escolas em 1992 não registraram crimes no \textit{campus}. Você acha que a falha das faculdades em informar os crimes pode ser vista como uma seleção amostral exógena? Explique.

\vspace{0.7cm}

\question\qstar (9.C14, \citep{wooldridge2016}) Use os dados do arquivo \texttt{ECONMATH} para responder a esta questão. O modelo populacional é
\[
score=\beta_0+\beta_1 act+u.
\]
\begin{enumerate}[label=(\roman*)]
    \item A nota do ACT está ausente para quantos estudantes? Qual é a fração da amostra? Defina uma nova variável, \textit{actmiss}, que é igual a um quando \textit{act} estiver ausente e zero quando não.
    \item Crie uma nova variável, digamos \textit{act0}, que é a nota do \textit{act} quando ela é relatada e zero quando \textit{act} estiver ausente. Encontre a média de \textit{act0} e compare-a com a média de \textit{act}.
    \item Faça a regressão simples de \textit{score} sobre \textit{act} usando somente os casos completos. O que você obtém para o coeficiente de inclinação e seu erro padrão robusto em relação à heteroscedasticidade?
    \item Faça a regressão simples de \textit{score} sobre \textit{act0} usando todos os casos. Compare o coeficiente de inclinação com o obtido no item (iii) e comente.
    \item Agora use todos os casos e faça a regressão
    \[
    score_i \text{ sobre } act0_i,\; actmiss_i.
    \]
    Qual é a estimativa de inclinação sobre \textit{act0}? De que forma isso se compara com suas respostas dos itens (iii) e (iv)?
    \item Comparando as regressões nos itens (iii) e (v), o uso de todos os casos e a adição de um estimador de dados ausentes melhora a estimação de \(\beta_1\)?
    \item Se você adicionar a variável \textit{colgpa} às regressões dos itens (iii) e (v), sua resposta ao item (vi) mudará?
\end{enumerate}

\end{questions}
\newpage
\bibliographystyle{apalike}
\bibliography{refs}
\end{document}
